\section{The Lupercalia of Mad Men}

\begin{quotex}
A time is coming when men will go mad, and when they see someone who is not mad, they will attack him, saying, ``You are mad, you are not like us.''
\flright{\textsc{St Anthony the Great}}
\end{quotex}


\paragraph{Guenon the Master}


\textbf{Rene Guenon} is the master who made religion and spirituality comprehensible to modern men. This was accomplished by pointing out the esoteric element, elucidating the true meaning of symbols, and by describing the metaphysical principles underlying Tradition. In the West, this metaphysic is most clearly expressed in the Platonic tradition, although, as we pointed out\footnote{\url{https://gornahoor.net/?p=1256}}, Plato was but one link in a much larger chain of Hermetic philosophy.

Now this manner of thinking used to be taught to educated men, but in our time, may seem quite incomprehensible. All the more so, it is worth the effort and the result is an intellectual conversion that will see the material as the effect of the spiritual, rather than the opposite.

Since \textbf{Julius Evola} is often the ``gateway drug'' to tradition, and since he made a conscious effort, in his philosophy of magical idealism, to become part of that chain of Hermetic philosophy, we take the trouble to make known the general outline of magical idealism. The ideal is to build on Evola, and even Guenon, and not fall into a form of sterile scholasticism.

\paragraph{What is Rationality}


\emph{Soi-disant} bright men today claim to be adherents of a scientific rationality, without really defining it. Rationality, if it means anything, is equivalent to knowing sufficient reason of things, i.e., why it is one thing rather than another. In that case science does not know the sufficient reason for the existence of the universe itself, life, or conscious thought. It never will, so we need to look elsewhere. As esoterists, we look within, since the truth resides in man's interiority (\textbf{St Augustine}).

As Evola points out, behind rationality, there is the alogon, which transcends all rational thought. Although he would not likely approve this formulation, it is still worth the effort. That is, the Father-God transcends the world. He then generates the Son, or Logos, through whom the world is created. This doctrine is critical. \emph{Pace} Arianism, if the Son is a creature, then the Father could have created Chaos in opposition to Logos. However, that is not possible since the Logos is one with the Father.

This is very important in practical terms. In Essence, the Logos is perfect. However, we don't experience that in Existence. But the reason for that cannot be a contrary transcendent principle that opposes us. Rather, as Evola points out, it indicates a privation, i.e., something we should have, but fail to bring about. The first alternative would mean that we are the playthings of other forces; the second indicates that we are lacking in this or that power, something that can be rectified.

\paragraph{Evolution from Apes}


I recently saw one of those famous Internet quizzes that pretend to show how well you have imbibed the consensus reality of the men who have gone mad. In in, you are expected to agree with the statement, ``Man evolved from the Ape.''

Of course, neo-Darwinism claims no such thing. The most they can claim is that at one point there was a population of some defunct ape species, following by a population of hominids. But they cannot show any sufficient reason for that population change. That is, there is no chain of events leading from one species to another. By its own principles, in place of any rational explanation, it postulates a series of random mutations that resulted in the hominid.

In other words, there is no ``inner man'' just waiting to arise from the ape. Rather there is a rupture, even by its own principles. That is why Evola can say that the ape represents a privation of the man struggling to manifest in Existence. This position is different from alternatives such as creationism or intelligent design; it need not even dispute the fossil record.

\paragraph{Three I's}


There are two dimensions of nature. The physical or hylic, represented qualitatively by the four elements and the psychic represented by the fifth element of life, irreducible to matter. The pneumatic, supernatural or spiritual element, is the third dimension that transcends the two natural dimensions. This is verified esoterically in the idea of the three I's. The physical body has an I, i.e., its own center of intelligence. The psychic body has the I of the personality. These are given naturally to every man, although few become aware of the two different I's, falsely believing instead there is an inner unity.

Awareness of the Real I, on the other hand, depends on making certain efforts. A man who believe he is one and free, will not make those efforts. Nevertheless, Traditional metaphysics becomes more comprehensible to those who make such efforts. It also constitutes the truth of such doctrines, beyond any rational argument, logical proof, or scientific experiment.


\paragraph{The Purge and Patricide}

I had the occasion to view the movie, The Purge\footnote{\url{https://www.imdb.com/title/tt2184339/}}, recently. The premise is that the USA has been reconstituted on different grounds with new founding fathers. They implemented a festival called the ``purge'', one night a year, during which any crime is permitted. This is sort of a latter day Lupercalia\footnote{\url{https://en.wikipedia.org/wiki/Lupercalia}}. The state clearly claims the power to declare that what was considered evil, can now be called ``good''.

The protagonist, a caring father, became prosperous from selling high end security systems, required for protection during the purges. The mother is the perfect wife, even cooking dinner while wearing a dress and high heels. On the night of the purge, the father arms the alarm system. In direct disobedience to the Father, the adolescent son lets a stranger, clearly not a resident of the neighborhood, into the house for sanctuary. Then marauders, who were stalking him as part of the purge, demand that the Father return the stranger to them, or risk having his home invaded and his family killed.

The Father then captures the stranger, ties him up, and prepares to hand him over to the marauders. The morality of the Father is to enforce the boundaries of his property and to protect the lives of his family members. Then the mother, along with the children, has a crisis of conscience, fells compassion for the man, and convinces the father not to release him outside. So the Father betrays his own principles and adopts the moral principles of women and children. Hence a hellish night ensues, born of their own choices.

Nevertheless, the destruction of his home and accepting the risk to his family, are insufficient to redeem him. He is killed as an act of symbolic patricide. This seems to empower the mother who suddenly becomes a strong character. When her neighbors arrive, not to help her, but rather to kill her out of covetousness, the stranger suddenly appears and rescues her and her family. Thus the mother, in alliance with the stranger, without her husband and contra her neighbors, is ready to create a new world.


\flrightit{Posted on 2014-07-10 by Cologero}

\begin{center}* * *\end{center}

\begin{footnotesize}
\begin{sffamily}
\texttt{Tom Blanchard on 2014-07-11 at 11:53 said:}

I have not seen ``The Purge'', although this is a powerful summary of its moral message. It is amazing how deeply the ``morality of the mother and children'' has penetrated modern western consciousness, to the point where it is always assumed that it is the logical conclusion of Christian teaching (even by the representatives of the Roman Church). One has difficulties even attempting to make a case for the Father's morality in our setting; even the most high-minded opponent would counter that the Father's ``noble sacrifice'' to die for this stranger's sake represents the pinnacle of Christian virtue and that to turn away the stranger would be an expression of imperfection and a short-sighted selfishness.

\hfill

\texttt{Max on 2014-07-12 at 10:25 said:}

When talking to people or looking around the internet, a general but vague awareness of something major changing is evident. People expect our ``civilization'' to burn, -we all know we deserve it. In fact it is what keeps some motivated to continue to live the lie. People want to experience the great collapse, to be in the middle of it. Secretly men wish for a Lupercalia, a scorched earth to renew the land. They say they will repopulate the earth and make a new order. Why would we believe them if they cannot produce something of that order in their own life right now? Guns and canned food will not make a new order.

Man seems destined to consume things increasingly more similar to his own nature. We made an alliance with animals to take care of and protect them and they gave us milk and wool. They trusted us but we betrayed them and put them in concentration camps. We made a prison for ourselves.

A problem arises if the scorchers only go half the way. Will the reality of the dissolution have them grasping for a false saviour and cling to every illusion just to stop their self induced suffering? Christ endured and he was resurrected, the true way of life.

The world is sacrificed to be resurrected. We cannot fight the collapse and we cannot further the collapse, only endure it and remember our beginnings with the knowledge that when beginning and end meet the world is united to eternity. It is only at the exact moment when the serpent bites its own tail that we can strike it with a single blow. Until that we will be here suffering. The fruit of knowledge of good and evil is transformed into the golden apple containing the liquid of immortality, blood of Christ. The cross leads us back to the tree of life, becoming visible again.

\hfill

\end{sffamily}
\end{footnotesize}

